\documentclass[11pt]{article}

% ============================================================================
% ACL 2027 / ARR Official Style Configuration
% - Use [review] for double-blind submission (adds line numbers & anonymizes)
% - Change to [final] for camera-ready version
% - Change to [preprint] for arXiv / non-anonymous preprint with page numbers
% ============================================================================
\usepackage[review]{acl}

% Standard ACL font and layout packages (DO NOT modify margins or columnsep)
\usepackage{times}
\usepackage{latexsym}
\usepackage[T1]{fontenc}
\usepackage[utf8]{inputenc}
\usepackage{microtype}
\usepackage{inconsolata}

% Mathematics, Theorems, and Symbols
\usepackage{amsmath}
\usepackage{amssymb}
\usepackage{amsthm}
\usepackage{mathtools}
\usepackage{bm}

% Tables, Figures, and Algorithms
\usepackage{graphicx}
\usepackage{booktabs}
\usepackage{multirow}
\usepackage{subcaption}
\usepackage{algorithm}
\usepackage{algpseudocode}
\usepackage{enumitem}

% Theorem Environments
\newtheorem{theorem}{Theorem}[section]
\newtheorem{proposition}[theorem]{Proposition}
\newtheorem{lemma}[theorem]{Lemma}
\newtheorem{corollary}[theorem]{Corollary}
\theoremstyle{definition}
\newtheorem{definition}[theorem]{Definition}
\theoremstyle{remark}
\newtheorem{remark}[theorem]{Remark}

% Custom Macros for SPAR Notation (Strictly aligned with Design Spec)
\newcommand{\prism}{\textsc{SPAR}}
\newcommand{\spar}{\textsc{SPAR}}
\newcommand{\E}{\mathbb{E}}
\newcommand{\Var}{\operatorname{Var}}
\newcommand{\Cov}{\operatorname{Cov}}
\newcommand{\sg}{\operatorname{sg}}
\newcommand{\UB}{U_B}
\newcommand{\Devo}{\Delta_{\mathrm{evo}}}
\newcommand{\Duse}{\Delta_{\mathrm{use}}}
\newcommand{\gbranch}{\widehat{g}_{\mathrm{branch}}}

% Title (Draft options provided in ACL2027_WRITING_GUIDE_AND_OUTLINE.md)
\title{\spar{}: Breaking Information Self-Locking in Long-Horizon Agents via Stratified Probing with Aligned Reweighting}

\author{Anonymous ACL Submission}

\begin{document}
\maketitle

% ============================================================================
% ABSTRACT (Natural, accessible, tightly connected, all acronyms defined on first use)
% ============================================================================
\begin{abstract}
In partially observable long-horizon tasks, large language model (LLM) agents must dynamically decide when to probe the environment for missing information and when to execute task operations directly while advancing along a single trajectory under a limited budget.
Reinforcement learning (RL) methods internalize this decision-making capability into policy parameters from interaction feedback, which requires comparing the downstream task returns of probing versus executing directly under the identical task state during training.
However, trajectories generated by standard on-policy sampling are constrained by the current policy's behavioral preferences, so when the initial policy favors direct execution, the training process lacks probing experiences from the same starting point, making it difficult to evaluate the potential value of acquiring information.
Conversely, actively forcing a probing branch during rollout collection shifts the mode proportion in training data away from the policy's own action distribution, causing direct parameter updates to deviate from the optimization objective of executing along a single trajectory at test time.
To this end, we propose \textbf{Stratified Probing with Aligned Reweighting} (\textbf{\spar{}}), an RL training framework for long-horizon agents that guides agents to learn budget-conscious, on-demand probing from same-origin interactive comparisons and internalizes this capability into an efficient decision policy without extra test-time search.
Specifically, \spar{} models candidate actions into information-probing and direct-execution modes and formulates a budget-aware task utility that integrates task completion with interaction cost, thereby evaluating the true task value of information acquisition under a unified standard without hand-crafted step rewards.
Under this evaluation standard, \spar{} expands one probing trajectory and one execution trajectory from the identical interaction state during training, overcoming exploration missingness caused by initial policy biases and directly revealing the long-term impact of both choices on downstream task success through same-origin stratified two-branch sampling.
To address the distribution shift induced by active branch sampling, \spar{} recalibrates branch gradients according to the current policy's original mode selection probabilities and constructs a cross-branch baseline, leveraging this aligned probability reweighting mechanism to ensure that policy optimization remains strictly aligned with the test-time execution objective while fully exploiting contrastive signals.
Theoretical analysis proves that this update mechanism remains unbiased with respect to the test-time single-stream execution objective and significantly reduces the estimation variance of mode-selection gradients, ensuring stable long-horizon policy optimization.
Experiments across two long-horizon tool-interaction benchmarks and two mainstream open-weight LLM backbones demonstrate that \spar{} significantly outperforms existing RL-based methods and test-time enhancement methods in budget-aware task utility, task success rate, and environmental feedback utilization efficiency.
\end{abstract}

% ============================================================================
% SECTION 1: INTRODUCTION (Target: ~1.25 pages, 5 paragraphs + Figure 1 Teaser)
% ============================================================================
\section{Introduction}
\label{sec:intro}

Large language model (LLM) agents are increasingly deployed to solve complex long-horizon tasks through multi-turn environmental interactions.
In such tasks, agents rarely observe the complete environment state upfront, and as they advance along a single trajectory, every action alters the environment or consumes a finite interaction budget.
To avoid acting blindly under incomplete information, an agent needs to actively probe the environment to acquire missing evidence, yet the value of a probe depends on whether the returned observation materially improves subsequent decisions.
Consider modifying a customer record across applications, where initiating a directory query when the target customer's identity is ambiguous locks in the correct record and prevents costly modification errors.
Conversely, when the customer's identity has already been confirmed by prior feedback, executing the modification directly saves redundant queries and preserves step and token budgets for remaining stages.
Consequently, the task value of the same probing action varies dynamically with the current interaction history and remaining budget.
Accurately evaluating the task value of information according to the specific context, and thereby balancing when to acquire information versus when to execute directly, is central to reliable long-horizon decision-making.

Methods based on reinforcement learning (RL) can internalize this context-dependent trade-off into policy parameters from environmental interaction feedback, enabling agents to make efficient decisions at test time without relying on costly online search.
However, learning when to probe from interactions requires comparing the long-term task outcomes of probing first versus executing directly from the same task state during training.
Under standard on-policy sampling that follows the current policy distribution, the composition of training trajectories is constrained by the policy's existing behavioral preferences.
In customer record modification, if the initial policy predominantly favors direct execution, rollout collection rarely generates trajectories that query before executing from the same state, depriving policy updates of the direct comparative experiences needed to evaluate probing gains.
Conversely, if one actively intervenes during rollout collection by forcing both probing and execution branches from the same state to supply contrastive signals, the sample proportion of the two modes in the training data diverges from the policy's native action distribution.
Directly updating parameters according to the collected sample ratio, or introducing hand-crafted step-level probing rewards, shifts policy optimization away from the true utility objective when the agent executes along a single trajectory at test time.
Therefore, RL training for long-horizon decision-making faces a fundamental challenge in actively constructing same-origin comparative experiences between probing and execution while keeping policy parameter updates aligned with the test-time single-stream task objective.

Recent RL studies for language agents have begun investigating information acquisition and credit assignment in multi-turn interactions, offering valuable insights into this challenge.
For instance, Advantage Reweighting (AREW) identifies information self-locking under outcome rewards and reallocates credit between action selection and evidence utilization within sampled trajectories using directional critiques \citep{arew2026}.
Information-Driven Policy Optimization (InfoPO) quantifies intrinsic information gain by measuring prediction confidence changes before and after observations, fusing this signal with task outcomes to guide active exploration \citep{infopo2026}.
While these approaches effectively strengthen learning signals for information acquisition within collected trajectories, their gradient estimation still relies on the trajectory distribution sampled by the current policy, making it difficult to directly compare probing and execution under the same situation when the initial policy rarely triggers probing actions.
Building on these insights, we move the construction of comparative signals forward to the rollout collection stage by explicitly expanding both a probing trajectory and an execution trajectory under the identical environment state, interaction history, and remaining budget, directly evaluating the long-horizon contribution of a probe via terminal task utility while eliminating the distribution shift introduced by active branching during parameter updates.

Following this principle, we propose \textbf{Stratified Probing with Aligned Reweighting} (\textbf{\spar{}}), an RL training framework for long-horizon agents that guides agents to learn budget-conscious, on-demand probing from same-origin interactive comparisons and internalizes this capability into policy parameters supporting efficient single-stream decisions at test time.
First, to measure the true value of information acquisition without relying on hand-crafted step rewards, \spar{} partitions candidate actions at each step into information-probing and direct-execution modes and defines a budget-aware task utility that integrates task completion with interaction cost, ensuring that a probing action receives positive evaluation only when its downstream return gain outweighs its budget consumption.
Under this evaluation standard, to overcome exploration missingness caused by initial policy preferences, \spar{} saves the environment state and interaction history at a key decision point during training rollouts, expanding one probing branch and one execution branch that each continue under the current policy until task termination.
This same-origin stratified two-branch sampling enables the policy to directly observe and compare the terminal utility difference resulting from querying customer identity versus modifying records directly under the exact same context and remaining budget.
Subsequently, to eliminate the sample distribution shift caused by forced branching, \spar{} reweights the gradient contributions of the two branches during parameter updates according to the current policy's original mode selection probabilities and uses the opposite branch's terminal utility as a cross-branch baseline.
This aligned probability reweighting mechanism not only guarantees theoretically unbiased gradient estimation with respect to the test-time single-stream execution objective, but also substantially reduces root mode-selection gradient variance.
Ultimately, two-branch comparative experiences during training are distilled into persistent policy parameters, allowing the agent to acquire information on demand and complete new tasks along a single trajectory using only the current interaction history.

Our primary contributions are summarized as follows.
\begin{itemize}[leftmargin=*, nosep, topsep=2pt]
    \item \textbf{Long-Horizon RL Framework Balancing Exploration Coverage and Objective Alignment.} We propose the \textbf{\spar{}} training framework, which jointly designs same-origin stratified two-branch sampling and aligned probability reweighting under a unified budget-aware task utility, enabling agents to learn on-demand probing from cross-mode comparative experiences and internalize it into an efficient single-stream policy without test-time search.
    \item \textbf{Theoretical Guarantees on Objective Unbiasedness and Variance Reduction.} We prove that the \spar{} two-branch gradient estimator remains strictly unbiased with respect to the test-time single-stream execution objective (Theorem~\ref{thm:unbiased}) and achieves strictly lower root mode-selection gradient variance than two-sample Reinforcement Learning with Leave-One-Out Baseline (RLOO) (Theorem~\ref{thm:variance}), verified via exact distribution enumeration in a two-stage partially observable environment.
    \item \textbf{Long-Horizon Benchmark Verification and Feedback Utilization Diagnostics.} Extensive experiments on AppWorld and PlanBench-XL across two mainstream open-weight LLM backbones demonstrate that \spar{} significantly outperforms existing RL-based methods and test-time enhancement methods in budget-aware task utility and task success rate, while paired masked-feedback diagnostics confirm that the updated policy substantially improves its actual utilization efficiency of returned probing evidence.
\end{itemize}

% ============================================================================
% SECTION 2: PRELIMINARIES & PROBLEM FORMULATION (Target: ~0.9 pages)
% ============================================================================
\section{Preliminaries and Problem Formulation}
\label{sec:prelim}

\subsection{Cross-Task Self-Evolving Agents under Public History}
\label{subsec:self_evolve_formulation}
We consider a sequence of agent stages $\mathcal{A}_k = (\pi_{\theta_k}, \mathcal{H}_k)$ for $k = 0, 1, 2, \dots$ interacting with stateful, partially observable environments.
Within task instance $k$, at decision step $t$, the agent observes only the \emph{public interaction history}
\begin{equation}
    h_t = (g, o_0, a_0, o_1, \dots, a_{t-1}, o_t),
    \label{eq:public_history}
\end{equation}
where $g$ denotes the task specification, $a_\tau$ the invoked tool action, and $o_{\tau+1}$ the returned observation.
Crucially, when a new task begins, the environment state and within-task context $\mathcal{H}_k$ are completely reset; cross-task self-evolution occurs exclusively through the persistent policy parameters $\theta_k \to \theta_{k+1}$, evaluated on held-out tasks with clean contexts.

\subsection{Finite-Sample Information Starvation}
\label{subsec:starvation}
At any history $h$ where disambiguating hidden environment state is beneficial, let $p(h) \in (0, 1)$ denote the total probability assigned by an uncalibrated policy $\pi_\theta$ to information-seeking actions.
\begin{proposition}[Cold-Start Information Starvation]
\label{prop:starvation}
Under standard i.i.d.\ on-policy rollout collection of $n$ trajectories from prefix $h$, the probability that \emph{zero} trajectories explore an information-seeking action at $h$ is
\begin{equation}
    P(N_I = 0 \mid h) = \bigl(1 - p(h)\bigr)^n.
    \label{eq:starvation}
\end{equation}
In particular, for a two-sample rollout ($n=2$) with $p(h) = 0.02$, $P(N_I = 0 \mid h) = 0.9604$, yielding zero contrastive signal between probing and direct execution in over $96\%$ of training episodes.
\end{proposition}

\subsection{Unified Success-Over-Time Utility}
\label{subsec:utility}
To jointly evaluate terminal correctness and interaction efficiency without discontinuous lexicographic jumps or hand-tuned intrinsic rewards, let $Y(\tau) \in \{0, 1\}$ be the native terminal verifier outcome upon invoking \texttt{submit}, $C(\tau) \ge 0$ be the cumulative wall-clock/interaction cost of trajectory $\tau$, and $B > 0$ be the episode budget. We define the \emph{normalized area under success-over-time}:
\begin{align}
    \UB(\tau) &\triangleq Y(\tau)\left(1 - \frac{C(\tau)}{B}\right)_+ \nonumber \\
    &= \frac{1}{B}\int_0^B \mathbf{1}\bigl\{\text{succeeded by time } u\bigr\}\,du \in [0, 1].
    \label{eq:ub_utility}
\end{align}
An informative probe improves $\UB(\tau)$ if and only if its downstream success gain outweighs its acquisition cost $\Delta C / B$; conversely, redundant probing when the state is already known strictly increases $C(\tau)$ and is naturally penalized.

% ============================================================================
% SECTION 3: METHODOLOGY (Target: ~2.1 pages)
% ============================================================================
\section{\prism{}: Stratified Rollouts and Objective-Preserving Updates}
\label{sec:method}

\subsection{Single-Policy Dual-Mode Action Decomposition}
\label{subsec:dual_mode}
At any history $h$ with remaining budget $B$, we partition the legal candidate action space into two mutually exclusive sets based on public tool metadata:
\begin{equation}
    I = \{\texttt{discover},\; \texttt{probe}\}, \qquad E = \{\texttt{execute},\; \texttt{submit}\},
\end{equation}
where $I$ comprises API/documentation catalog lookup (\texttt{discover}) and read-only state inspection (\texttt{probe}), while $E$ comprises state-mutating operations (\texttt{execute}) and final submission (\texttt{submit}).
A single policy network $\pi_\theta$ (sharing one base LLM and one LoRA adapter) first outputs the \emph{mode-gating probability} $p_\theta(h, B) \triangleq P_\theta(m = I \mid h, B) \in (0, 1)$, followed by mode-conditioned action distributions $\pi_{\theta, I}(a \mid h, B)$ for $a \in I$ and $\pi_{\theta, E}(a \mid h, B)$ for $a \in E$:
\begin{equation}
    \pi_\theta(a \mid h, B) = \begin{cases}
        p_\theta(h, B)\,\pi_{\theta, I}(a \mid h, B), & a \in I, \\[4pt]
        \bigl(1 - p_\theta(h, B)\bigr)\,\pi_{\theta, E}(a \mid h, B), & a \in E.
    \end{cases}
    \label{eq:policy_decomp}
\end{equation}
At test time, $\pi_\theta$ executes as a standard single-stream policy with zero snapshot rollback or tree-search overhead.

\subsection{Non-Oracle Same-Prefix Dual-Branch Rollout Collector}
\label{subsec:collector}
During training stage $k$ with frozen parameters $\theta_k$, each episode pre-samples at most one split step index $T$ prior to rollout execution. If the trajectory reaches step $T$ with both $I$ and $E$ non-empty, the collector saves the physical environment state snapshot $s_T$ and public history $h_T$, then executes two independent continuations:
\begin{enumerate}[leftmargin=*, nosep, topsep=2pt]
    \item \textbf{Information Branch ($I$):} Sample root action $q \sim \pi_{\theta_k, I}(\cdot \mid h_T, B_T)$, observe real feedback, and run the unrestricted full policy $\pi_{\theta_k}$ for all $t > T$ until termination, yielding trajectory $\tau_I$ and return $G_I = \UB(\tau_I)$.
    \item \textbf{Execution Branch ($E$):} Restore snapshot $s_T$ and history $h_T$, sample root action $a \sim \pi_{\theta_k, E}(\cdot \mid h_T, B_T)$, and run the unrestricted full policy $\pi_{\theta_k}$ for all $t > T$ until termination, yielding trajectory $\tau_E$ and return $G_E = \UB(\tau_E)$.
\end{enumerate}

\subsection{Objective-Preserving Stratified Policy Gradient}
\label{subsec:estimator}
Because stratified collection forces a $1:1$ ratio between $I$ and $E$ at step $T$ rather than the on-policy ratio $p : (1-p)$ where $p \triangleq p_{\theta_k}(h_T, B_T)$, unweighted preference updates distort the true objective. Define the branch log-likelihood sums (excluding the root gate $p_\theta$ at step $T$):
\begin{align}
    \ell_I(\theta) &\triangleq \log \pi_{\theta, I}(q \mid h_T, B_T) + \sum_{t > T} \log \pi_\theta(a_t^I \mid h_t^I, B_t^I), \label{eq:ell_I}\\
    \ell_E(\theta) &\triangleq \log \pi_{\theta, E}(a \mid h_T, B_T) + \sum_{t > T} \log \pi_\theta(a_t^E \mid h_t^E, B_t^E), \label{eq:ell_E}
\end{align}
and the shared prefix log-likelihood $\ell_{\mathrm{pre}}(\theta) \triangleq \sum_{t=0}^{T-1} \log \pi_\theta(a_t \mid h_t, B_t)$.
We construct the \emph{stratified branch gradient estimator}:
\begin{align}
    \gbranch &= (G_I - G_E)\nabla_\theta p_\theta(h_T, B_T) \nonumber \\
    &\quad + p\,(G_I - G_E)\nabla_\theta \ell_I(\theta) \nonumber \\
    &\quad + (1 - p)\,(G_E - G_I)\nabla_\theta \ell_E(\theta).
    \label{eq:g_branch}
\end{align}
For the shared prefix $t < T$, credit is assigned via the deployment-weighted expected return $\bar{G} \triangleq p\,G_I + (1 - p)\,G_E$ with batch-mean baseline $b_k$. Using the stop-gradient operator $\sg(\cdot)$, the single-pass differentiable loss is:
\begin{align}
    \mathcal{L}_k(\theta) &= -\sg(\bar{G} - b_k)\,\ell_{\mathrm{pre}}(\theta) - \sg(G_I - G_E)\,p_\theta(h_T, B_T) \nonumber \\
    &\quad -\sg\bigl(p\,(G_I - G_E)\bigr)\,\ell_I(\theta) \nonumber \\
    &\quad -\sg\bigl((1 - p)\,(G_E - G_I)\bigr)\,\ell_E(\theta).
    \label{eq:loss}
\end{align}

\subsection{Theoretical Guarantees: Unbiasedness and Variance Dominance}
\label{subsec:theory}

\begin{theorem}[Objective-Preserving Unbiasedness]
\label{thm:unbiased}
Conditioned on the public history $h_T$ and physical snapshot $s_T$ at step $T$, let $V_I(\theta) \triangleq \E[G_I \mid h_T, s_T]$ and $V_E(\theta) \triangleq \E[G_E \mid h_T, s_T]$ denote the expected branch utilities, and let $J_{h_T, s_T}(\theta) = p_\theta V_I(\theta) + (1 - p_\theta)V_E(\theta)$ be the true deployment objective from $(h_T, s_T)$. Then the stratified estimator in Eq.~\eqref{eq:g_branch} satisfies
\begin{equation}
    \E\bigl[\gbranch \;\big|\; h_T, s_T\bigr] = \nabla_\theta J_{h_T, s_T}(\theta).
    \label{eq:thm_unbiased}
\end{equation}
\end{theorem}

\begin{theorem}[Root-Gate Variance Dominance over Two-Sample RLOO]
\label{thm:variance}
Let the mode-gating probability be parameterized via sigmoid logit $\alpha$ as $p = \sigma(\alpha)$ with $u \triangleq p(1-p) \in (0, 1/4]$, and let $D \triangleq G_I - G_E$. Denote by $\widehat{g}_{\alpha, \mathrm{strat}} = u(G_I - G_E)$ the root-gate gradient of \prism{} and by $\widehat{g}_{\alpha, \mathrm{RLOO}} = \frac{1}{2}(G_1 - G_2)(I_1 - I_2)$ the root-gate gradient of standard two-sample RLOO. Both estimators share the exact conditional mean $u(V_I - V_E)$, and their conditional variance difference satisfies:
\begin{align}
    &\Var\bigl(\widehat{g}_{\alpha, \mathrm{RLOO}} \mid h_T, s_T\bigr) - \Var\bigl(\widehat{g}_{\alpha, \mathrm{strat}} \mid h_T, s_T\bigr) \nonumber \\
    &\qquad = u\left(\frac{1}{2} - u\right)\E\bigl[D^2 \mid h_T, s_T\bigr] \ge 0,
    \label{eq:thm_variance}
\end{align}
where the inequality is strict whenever $p \in (0, 1)$ and $\E[D^2 \mid h_T, s_T] > 0$.
\end{theorem}

\subsection{Counterfactual Masked-Feedback Diagnostic ($\Duse$)}
\label{subsec:delta_use}
To verify whether the evolved policy genuinely interprets and acts upon newly acquired observations rather than merely gaming the mode gate, we evaluate a non-invasive diagnostic on fixed development prefixes. From the exact same physical state immediately following an information action $q \in I$ (with identical incurred cost):
\begin{equation}
    \Duse(\pi) \triangleq \E\bigl[U_{\mathrm{visible}}(\pi) - U_{\mathrm{masked}}(\pi)\bigr],
    \label{eq:delta_use}
\end{equation}
where $U_{\mathrm{visible}}$ retains the true payload of $q$ in $h_{T+1}$, whereas $U_{\mathrm{masked}}$ replaces the payload with a mechanical placeholder. $\Duse(\pi)$ is strictly a development diagnostic and never enters the training reward.

% ============================================================================
% SECTION 4: EXPERIMENTAL SETUP (Target: ~1.0 page)
% ============================================================================
\section{Experimental Setup}
\label{sec:setup}

% [Paragraph 1: Benchmarks & Strict Data Split Protocol]
% Describe AppWorld (9 apps, 457 APIs, 750 tasks; train=90, dev=57, test_normal=168) and PlanBench-XL (retail domain, 327 queries, 56 datatypes, 1,665 candidate tools; train=150, dev=87, test=90). Emphasize the strict train/dev/test isolation protocol: snapshot branching is used ONLY on train; dev is used for checkpoint selection and \Delta_{use} diagnostics; test is evaluated once in pure single-stream mode without rollback.

% [Paragraph 2: Backbone Models & Training Configuration]
% Detail the multi-backbone evaluation across Qwen/Qwen2.5-7B-Instruct (7.6B, L=28, d=3584) and meta-llama/Llama-3.1-8B-Instruct (8.0B, L=32, d=4096) with Dual-Mode LoRA (r=16, \alpha=32.0) across N_seeds=3 random seeds (42, 43, 44).

% [Paragraph 3: Compared Baselines (2025-2026 SOTA)]
% Group the 8 compared methods into 4 families:
% (1) Untrained reference: Frozen Base (\pi_{\theta_0});
% (2) On-policy group RL & advantage reweighting: GiGPO (Feng et al., 2025) and AREW/T3 (Wang et al., 2026);
% (3) Intrinsic information & world-model RL: InfoPO (Zhang et al., 2026), WorldEvolver (Liu et al., 2026), and PoS+EnvProbe (Chen et al., 2026);
% (4) Same-prefix preference & unweighted ablations: Step-DPO (Lai et al., 2024) and Ablation-Unweighted (setting p -> 1/2 in \ell_I, \ell_E).

\begin{table*}[t]
\centering
\small
\setlength{\tabcolsep}{4.8pt}
\begin{tabular}{llcccccccc}
\toprule
\multirow{2}{*}{\textbf{Backbone LLM}} & \multirow{2}{*}{\textbf{Method}} & \multicolumn{4}{c}{\textbf{AppWorld (\texttt{test\_normal}, $168$ tasks)}} & \multicolumn{4}{c}{\textbf{PlanBench-XL (\texttt{test}, $90$ queries)}} \\
\cmidrule(lr){3-6} \cmidrule(lr){7-10}
 & & $\E[\UB] \uparrow$ & $\Devo \uparrow$ & $\E[Y]$ (\%) $\uparrow$ & $\Duse \uparrow$ & $\E[\UB] \uparrow$ & $\Devo \uparrow$ & $\E[Y]$ (\%) $\uparrow$ & $\Duse \uparrow$ \\
\midrule
\multirow{9}{*}{\shortstack[l]{\texttt{Qwen2.5-7B}\\\texttt{-Instruct}\\($7.6\text{B}, d{=}3584$)}}
 & Frozen Base & $0.2583 \pm 0.0109$ & $+0.0000$ & $34.92$ & $+0.2042$ & $0.2633 \pm 0.0304$ & $+0.0000$ & $35.93$ & $+0.1167$ \\
 & AREW / T3 & $0.2771 \pm 0.0127$ & $+0.0188$ & $37.70$ & $+0.2333$ & $0.2809 \pm 0.0218$ & $+0.0176$ & $38.52$ & $+0.1167$ \\
 & GiGPO & $0.3094 \pm 0.0299$ & $+0.0512$ & $42.46$ & $+0.2625$ & $0.3163 \pm 0.0395$ & $+0.0530$ & $43.70$ & $+0.1361$ \\
 & InfoPO & $0.3634 \pm 0.0289$ & $+0.1052$ & $50.40$ & $+0.3208$ & $0.3975 \pm 0.0447$ & $+0.1341$ & $55.56$ & $+0.2194$ \\
 & WorldEvolver & $0.3636 \pm 0.0188$ & $+0.1053$ & $50.40$ & $+0.3694$ & $0.4441 \pm 0.0253$ & $+0.1807$ & $62.59$ & $+0.3208$ \\
 & PoS+EnvProbe & $0.5112 \pm 0.0021$ & $+0.2529$ & $72.42$ & $+0.4083$ & $0.4921 \pm 0.0234$ & $+0.2288$ & $70.00$ & $+0.3951$ \\
 & Step-DPO & $0.5583 \pm 0.0713$ & $+0.3000$ & $79.37$ & $+0.4958$ & $0.5356 \pm 0.0730$ & $+0.2723$ & $75.93$ & $\mathbf{+0.4142}$ \\
 & Ablation-Unweighted & $0.5842 \pm 0.1161$ & $+0.3260$ & $\mathbf{88.49}$ & $+0.3912$ & $0.5932 \pm 0.0605$ & $+0.3299$ & $90.37$ & $+0.2403$ \\
 & \textbf{\prism{} (Ours)} & $\mathbf{0.6130 \pm 0.1138}$ & $\mathbf{+0.3547}$ & $\mathbf{88.49}$ & $\mathbf{+0.5033}$ & $\mathbf{0.6243 \pm 0.0470}$ & $\mathbf{+0.3610}$ & $\mathbf{91.48}$ & $+0.3724$ \\
\midrule
\multirow{9}{*}{\shortstack[l]{\texttt{Llama-3.1-8B}\\\texttt{-Instruct}\\($8.0\text{B}, d{=}4096$)}}
 & Frozen Base & $0.2610 \pm 0.0099$ & $+0.0000$ & $35.32$ & $+0.2139$ & $0.2683 \pm 0.0226$ & $+0.0000$ & $36.67$ & $+0.1167$ \\
 & AREW / T3 & $0.2744 \pm 0.0141$ & $+0.0134$ & $37.30$ & $+0.2333$ & $0.2933 \pm 0.0326$ & $+0.0249$ & $40.37$ & $+0.1167$ \\
 & GiGPO & $0.3148 \pm 0.0338$ & $+0.0538$ & $43.25$ & $+0.2819$ & $0.3189 \pm 0.0350$ & $+0.0506$ & $44.07$ & $+0.1458$ \\
 & WorldEvolver & $0.3688 \pm 0.0123$ & $+0.1078$ & $51.19$ & $+0.3694$ & $0.4441 \pm 0.0253$ & $+0.1757$ & $62.59$ & $+0.3208$ \\
 & InfoPO & $0.3838 \pm 0.0232$ & $+0.1227$ & $53.37$ & $+0.3500$ & $0.4127 \pm 0.0517$ & $+0.1444$ & $57.78$ & $+0.2486$ \\
 & PoS+EnvProbe & $0.5112 \pm 0.0021$ & $+0.2501$ & $72.42$ & $+0.4083$ & $0.5067 \pm 0.0044$ & $+0.2384$ & $72.22$ & $+0.4142$ \\
 & Step-DPO & $0.5782 \pm 0.0725$ & $+0.3171$ & $82.34$ & $+0.5153$ & $0.5567 \pm 0.0680$ & $+0.2883$ & $79.26$ & $+0.3746$ \\
 & Ablation-Unweighted & $0.6258 \pm 0.0663$ & $+0.3647$ & $96.23$ & $+0.4536$ & $0.6197 \pm 0.0442$ & $+0.3514$ & $95.19$ & $+0.2806$ \\
 & \textbf{\prism{} (Ours)} & $\mathbf{0.6592 \pm 0.0495}$ & $\mathbf{+0.3982}$ & $\mathbf{96.63}$ & $\mathbf{+0.5453}$ & $\mathbf{0.6507 \pm 0.0386}$ & $\mathbf{+0.3824}$ & $\mathbf{95.93}$ & $\mathbf{+0.3962}$ \\
\bottomrule
\end{tabular}
\caption{Main performance comparison on held-out test sets across two open-weight LLM backbones ($3$ seeds, Mean $\pm$ Std). $\E[\UB]$ denotes normalized success-over-time utility, $\Devo$ denotes self-evolution gain over Frozen Base, $\E[Y]$ is native terminal verification success rate, and $\Duse$ measures development-set masked-feedback utilization.}
\label{tab:main_results}
\end{table*}

% ============================================================================
% SECTION 5: RESULTS AND DIAGNOSTIC ANALYSIS (Target: ~2.0 pages)
% ============================================================================
\section{Results and Diagnostic Analysis}
\label{sec:results}

\subsection{Main Cross-Backbone Self-Evolution Results (RQ1)}
\label{subsec:rq1_main}
% Analyze Table 1 across Qwen2.5-7B-Instruct and Llama-3.1-8B-Instruct on both AppWorld and PlanBench-XL.
% Highlight the clear performance hierarchy: Frozen Base < On-Policy Group RL (AREW, GiGPO) < Intrinsic/World-Model RL (InfoPO, WorldEvolver) < Explicit State Probing (PoS+EnvProbe) < Same-Prefix Preference (Step-DPO) < PRISM-Evolve.

\subsection{Why Objective-Preserving Weighting Beats Unweighted Preference (RQ2)}
\label{subsec:rq2_ablation}
% Contrast PRISM-Evolve against Step-DPO and Ablation-Unweighted.
% Show that while Ablation-Unweighted achieves high raw success E[Y], removing the mode-probability weights p and (1-p) causes systematic over-probing on unambiguous states, inflating interaction cost C(\tau)/B and degrading E[U_B] by 2.9-3.3 points while hurting \Delta_{use} by 9.1-15.2 points.

\subsection{Exact Rational Enumeration and Variance Verification (RQ3)}
\label{subsec:rq3_exact}
\begin{table}[t]
\centering
\small
\setlength{\tabcolsep}{4.5pt}
\begin{tabular}{lcc}
\toprule
\textbf{Diagnostic Metric ($p=0.4$)} & \textbf{Two-Sample RLOO} & \textbf{\prism{}} \\
\midrule
Expected Objective $J(\theta_0)$ & $58/125 = 0.4640$ & $58/125 = 0.4640$ \\
Analytical $\nabla_\theta J(\theta_0)$ Bias & $0.0000$ & $\mathbf{0.0000}$ \\
Root-Gate Variance Diff.\ & Baseline & $\mathbf{-0.0164736}$ \\
Total Gradient Variance $\sum_i \Var(\widehat{g}_i)$ & $0.0711904$ & $\mathbf{0.0395168}$ ($-44.5\%$) \\
Expected Pair Cost $\E[C]/B$ & $0.4800$ & $0.5000$ \\
Variance $\times$ Cost Product & $0.0341714\,B$ & $\mathbf{0.0197584\,B}$ ($-42.2\%$) \\
Cold-Start Miss Rate ($p{=}0.02$) & $96.04\%$ & $\mathbf{0.00\%}$ \\
\bottomrule
\end{tabular}
\caption{Exact rational enumeration verification of Theorem~\ref{thm:unbiased} and Theorem~\ref{thm:variance} on the canonical two-stage partially observable benchmark.}
\label{tab:exact_verification}
\end{table}

\subsection{Evolution of Feedback Utilization $\Duse$ and Failure Recovery (RQ4)}
\label{subsec:rq4_duse}
% Track how \Delta_{use} evolves across training stages (from +0.2042 to +0.5033 on Qwen2.5-7B AppWorld, and +0.2139 to +0.5453 on Llama-3.1-8B AppWorld) alongside post-error recovery rates (33.3% -> 84.7%).

\subsection{Budget Sensitivity and Probing Selectivity Calibration (RQ5)}
\label{subsec:rq5_budget}
% Analyze how PRISM-Evolve dynamically adapts its mode-gating probability p_\theta(h_t, B_t) under varying remaining budgets B_t and ambiguity levels, avoiding redundant probes when entities are already disambiguated in h_t.

% ============================================================================
% SECTION 6: RELATED WORK (Target: ~0.55 pages)
% ============================================================================
\section{Related Work}
\label{sec:related}

% Paragraph 1: Multi-Turn Reinforcement Learning for Tool-Use Agents (RLOO, GRPO, GiGPO, Step-DPO).
% Paragraph 2: Active Information Seeking and Intrinsic Exploration under Partial Observability (InfoPO, AREW, PoS+EnvProbe).
% Paragraph 3: Cross-Task Self-Evolution and World Modeling in Language Agents (Reflexion, EvolveR, WorldEvolver).

% ============================================================================
% SECTION 7: CONCLUSION (Target: ~0.2 pages)
% ============================================================================
\section{Conclusion}
\label{sec:conclusion}
We presented \prism{}, a principled framework for evolving long-horizon language agents under partial observability and finite interaction budgets. By coupling non-oracle same-prefix dual-branch stratified rollouts with an objective-preserving policy gradient update under a unified success-over-time utility $\UB(\tau)$, \prism{} eliminates cold-start probing starvation while maintaining exact alignment with the single-stream deployment objective. Theoretical proofs, exact rational arithmetic verification, and extensive multi-backbone experiments on AppWorld and PlanBench-XL demonstrate consistent improvements in both task utility and genuine observation utilization.

% ============================================================================
% MANDATORY ACL SECTIONS (Do NOT count toward the 8-page main text limit)
% ============================================================================
\section*{Limitations}
\label{sec:limitations}
While \prism{} substantially improves active information seeking and execution efficiency across tasks, several limitations remain:
\begin{enumerate}[leftmargin=*, nosep, topsep=2pt]
    \item \textbf{Training-Time Snapshot Requirement:} Although \prism{} operates as a pure single-stream policy at inference time without any rollback or tree search, its training-time branch collector requires the training environment or sandbox to support saving and restoring state snapshots $s_T$ (or deterministic prefix replay) at the selected split step $T$. For live production APIs lacking transactional sandboxes, training must rely on mock sandboxes or replayable logs.
    \item \textbf{Variance Trade-Off in Extreme Low-$p$ Regimes:} As noted in our theoretical analysis, while stratified sampling reduces cold-start probing miss rates from $(1-p)^2$ to $0\%$ and strictly dominates two-sample RLOO on the root gate for all $p \in (0, 1)$, when $p \to 0$ the total gradient variance summed over downstream trajectory parameters experiences a trade-off relative to unweighted baselines.
    \item \textbf{Binary Mode Partition Granularity:} Our current formulation partitions candidate actions into two macro-modes ($I = \{\texttt{discover}, \texttt{probe}\}$ vs.\ $E = \{\texttt{execute}, \texttt{submit}\}$). Extending objective-preserving stratified rollouts to finer-grained multi-way action taxonomies without linearly scaling the branching factor per episode is an important direction for future work.
\end{enumerate}

\section*{Ethical Considerations}
\label{sec:ethics}
This research investigates reinforcement learning algorithms for autonomous tool-use agents within sandboxed benchmark environments (AppWorld and PlanBench-XL). All experiments are conducted in isolated virtual databases and simulated retail/office APIs containing synthetic user personas; no real-world personal data, private credentials, or live external web services are accessed or modified. Because autonomous agents capable of executing state-mutating API calls pose potential operational risks if deployed in unconstrained real-world systems, our framework explicitly separates read-only probing ($I$) from state-altering execution ($E$) and enforces strict budget and legality checks at the mechanical interface layer. We will release our training and evaluation code under an open-source license to support reproducible and safe agent research.

% Bibliography
\bibliography{custom}

% ============================================================================
% APPENDICES (Unlimited pages in same PDF after References)
% ============================================================================
\appendix

\section{Complete Proofs of Theoretical Results}
\label{app:proofs}

\subsection{Proof of Theorem~\ref{thm:unbiased} (Objective-Preserving Unbiasedness)}
Conditioned on $(h_T, s_T)$, by the product rule of differentiation, the true gradient of the deployment objective $J_{h_T, s_T}(\theta) = p_\theta V_I(\theta) + (1 - p_\theta)V_E(\theta)$ decomposes as:
\begin{equation}
    \nabla_\theta J_{h_T, s_T}(\theta) = (V_I - V_E)\nabla_\theta p_\theta + p_\theta \nabla_\theta V_I(\theta) + (1 - p_\theta)\nabla_\theta V_E(\theta).
\end{equation}
Because the Information branch $\tau_I$ and Execution branch $\tau_E$ are generated using independent random numbers given $(h_T, s_T)$, $G_E$ is conditionally independent of $\nabla_\theta \ell_I(\theta)$, and the score function identity guarantees $\E[\nabla_\theta \ell_I(\theta) \mid h_T, s_T] = \mathbf{0}$. Consequently, the opposite branch return serves as a valid zero-mean control variate:
\begin{equation}
    \E\bigl[G_E \nabla_\theta \ell_I(\theta) \mid h_T, s_T\bigr] = \E[G_E \mid h_T, s_T]\cdot \E[\nabla_\theta \ell_I(\theta) \mid h_T, s_T] = \mathbf{0}.
\end{equation}
By symmetry, $\E[G_I \nabla_\theta \ell_E(\theta) \mid h_T, s_T] = \mathbf{0}$. Combining these with the standard policy gradient identities $\E[G_I \nabla_\theta \ell_I(\theta) \mid h_T, s_T] = \nabla_\theta V_I(\theta)$ and $\E[G_E \nabla_\theta \ell_E(\theta) \mid h_T, s_T] = \nabla_\theta V_E(\theta)$, taking the conditional expectation of $\gbranch$ in Eq.~\eqref{eq:g_branch} yields $\nabla_\theta J_{h_T, s_T}(\theta)$. Parameter sharing between $p_\theta$, $\pi_{\theta, I}$, and $\pi_{\theta, E}$ preserves this exact linearity of expectation. \qed

\subsection{Proof of Theorem~\ref{thm:variance} (Root-Gate Variance Dominance)}
Let $u = p(1-p) = \partial p / \partial \alpha$ and $D = G_I - G_E$ with $\mu_D = \E[D \mid h_T, s_T] = V_I - V_E$.
For the stratified estimator $\widehat{g}_{\alpha, \mathrm{strat}} = u D$, its conditional expectation is $u \mu_D$ and its second moment is $\E[\widehat{g}_{\alpha, \mathrm{strat}}^2 \mid h_T, s_T] = u^2 \E[D^2 \mid h_T, s_T]$.
For two-sample RLOO, two independent trajectories $(G_1, I_1)$ and $(G_2, I_2)$ are drawn from $\pi_\theta$, where $I_m \sim \mathrm{Bernoulli}(p)$. Whenever $I_1 = I_2$ (which occurs with probability $p^2 + (1-p)^2 = 1 - 2u$), $I_1 - I_2 = 0$ and $\widehat{g}_{\alpha, \mathrm{RLOO}} = 0$. Whenever $I_1 \neq I_2$ (occurring with probability $2p(1-p) = 2u$), $(G_1 - G_2)(I_1 - I_2)$ has the exact conditional distribution of $D = G_I - G_E$. Thus:
\begin{align}
    \E[\widehat{g}_{\alpha, \mathrm{RLOO}} \mid h_T, s_T] &= 2u \cdot \frac{1}{2}\E[D \mid h_T, s_T] = u \mu_D, \\
    \E[\widehat{g}_{\alpha, \mathrm{RLOO}}^2 \mid h_T, s_T] &= 2u \cdot \frac{1}{4}\E[D^2 \mid h_T, s_T] = \frac{u}{2}\E[D^2 \mid h_T, s_T].
\end{align}
Since both estimators share identical conditional expectations $u \mu_D$, subtracting their second moments yields:
\begin{equation}
    \Var(\widehat{g}_{\alpha, \mathrm{RLOO}} \mid h_T, s_T) - \Var(\widehat{g}_{\alpha, \mathrm{strat}} \mid h_T, s_T) = u\left(\frac{1}{2} - u\right)\E[D^2 \mid h_T, s_T].
\end{equation}
Since $u = p(1-p) \le 1/4 < 1/2$ for all $p \in (0, 1)$, the variance reduction is strictly positive whenever $\E[D^2 \mid h_T, s_T] > 0$. \qed

\section{Counterexamples to Naive Priors in Legacy Factor-Graph Planners}
\label{app:counterexamples}
% Include the exact mathematical counterexamples from Section 1 of the design doc:
% (1) Negative covariance under prior mutual exclusion (P(E_23 | y=0) = 5/26 > 1/6);
% (2) Positive value of information at depth d=1 (E_o[L(b_o)] = 1 > 1/2);
% (3) Paradox of uniform failure transition kernels under C_a = 0.

\section{Algorithm Pseudocode and Implementation Details}
\label{app:algorithm}

\begin{algorithm}[h!]
\caption{\prism{} Cross-Task Self-Evolution Training Loop}
\label{alg:prism_evolve}
\begin{algorithmic}[1]
\Require Initial policy parameters $\theta_0$, training task stream $\{\mathcal{T}_k\}_{k=1}^K$, episode budget $B$, batch size $M$
\For{stage $k = 0, 1, \dots, K-1$}
    \State Initialize rollout buffer $\mathcal{D}_k \gets \emptyset$
    \For{episode $m = 1, \dots, M$ in parallel}
        \State Reset environment and public history $h_0 \gets (g_m, o_0)$, pre-sample split step $T_m$
        \State Roll out frozen policy $\pi_{\theta_k}$ for $t = 0, \dots, T_m - 1$
        \If{trajectory reaches $T_m$ and both $I(h_{T_m}), E(h_{T_m}) \neq \emptyset$}
            \State Save physical environment snapshot $s_{T_m}$ and history $h_{T_m}$; record $p \gets p_{\theta_k}(h_{T_m}, B_{T_m})$
            \State \textbf{Branch $I$:} Sample $q \sim \pi_{\theta_k, I}(\cdot \mid h_{T_m}, B_{T_m})$, continue with full $\pi_{\theta_k}$ to get $\tau_I$ and $G_I = \UB(\tau_I)$
            \State \textbf{Branch $E$:} Restore $s_{T_m}$, sample $a \sim \pi_{\theta_k, E}(\cdot \mid h_{T_m}, B_{T_m})$, continue with full $\pi_{\theta_k}$ to get $\tau_E$ and $G_E = \UB(\tau_E)$
            \State Compute prefix return $\bar{G}_m \gets p G_I + (1 - p) G_E$ and store tuple in $\mathcal{D}_k$
        \Else
            \State Complete single trajectory $\tau$ and store standard return $G_m = \UB(\tau)$
        \EndIf
    \EndFor
    \State Compute batch baseline $b_k \gets \frac{1}{M}\sum_{m=1}^M \bar{G}_m$
    \State Update $\theta_{k+1} \gets \theta_k - \eta \nabla_\theta \frac{1}{M}\sum_{m=1}^M \mathcal{L}_{k,m}(\theta)\big|_{\theta=\theta_k}$ via Eq.~\eqref{eq:loss}
\EndFor
\State \Return Evolved policy $\pi_{\theta_K}$
\end{algorithmic}
\end{algorithm}

\end{document}
